\documentclass{ximera}

\input{../../../preamble.tex}


\definecolor{fvdnavy}{HTML}{071D43}
\definecolor{fvdcyan}{HTML}{20BFC6}
\definecolor{fvdcyanlight}{HTML}{EFFCFB}
\definecolor{fvdmagenta}{HTML}{F10D69}
\definecolor{fvdmagentalight}{HTML}{FFF1F7}
\definecolor{fvdtext}{HTML}{163044}





%=================================================
% Pedagogische omgevingen
% PDF: tcolorbox
% HTML: learning-box
%=================================================

\newenvironment{voorkennis}
{%
    \ifdefined\HCode
        \HCode{%
<section class="learning-box learning-box--prior">
<div class="learning-box__header">
<span class="learning-box__icon" aria-hidden="true"></span>
<span class="learning-box__title">Voorkennis</span>
</div>
<div class="learning-box__content">
        }%
        \begin{itemize}
    \else
       \begin{tcolorbox}[
    enhanced,
    breakable,
    colback=fvdcyanlight,
    colframe=fvdcyan,
    coltext=fvdtext,
    boxrule=0.5pt,
    leftrule=4pt,
    arc=4mm,
    outer arc=4mm,
    left=7mm,
    right=6mm,
    top=5mm,
    bottom=5mm,
    before skip=6mm,
    after skip=6mm,
    title=Voorkennis,
    coltitle=fvdcyan,
    fonttitle=\bfseries\Large,
    attach title to upper={\par\smallskip},
    boxed title style={
        empty,
        left=0mm,
        right=0mm,
        top=0mm,
        bottom=0mm
    },
    overlay={
        \draw[
            fvdcyan,
            line width=0.5pt
        ]
        ([xshift=42mm,yshift=-13mm]frame.north west)
        --
        ([xshift=-7mm,yshift=-13mm]frame.north east);
    }
]
        \begin{itemize}
    \fi
}
{%
    \end{itemize}
    \ifdefined\HCode
        \HCode{%
</div>
</section>
        }%
    \else
        \end{tcolorbox}
    \fi
}


\newenvironment{leerdoelen}
{%
    \ifdefined\HCode
        \HCode{%
<section class="learning-box learning-box--goals">
<div class="learning-box__header">
<span class="learning-box__icon" aria-hidden="true"></span>
<span class="learning-box__title">Leerdoelen</span>
</div>
<div class="learning-box__content">
        }%
        \begin{itemize}
    \else
\begin{tcolorbox}[
    enhanced,
    breakable,
    colback=fvdmagentalight,
    colframe=fvdmagenta,
    coltext=fvdtext,
    boxrule=0.5pt,
    leftrule=4pt,
    arc=4mm,
    outer arc=4mm,
    left=7mm,
    right=6mm,
    top=5mm,
    bottom=5mm,
    before skip=6mm,
    after skip=6mm,
    title=Leerdoelen,
    coltitle=fvdmagenta,
    fonttitle=\bfseries\Large,
    attach title to upper={\par\smallskip},
    boxed title style={
        empty,
        left=0mm,
        right=0mm,
        top=0mm,
        bottom=0mm
    },
    overlay={
        \draw[
            fvdmagenta,
            line width=0.5pt
        ]
        ([xshift=42mm,yshift=-13mm]frame.north west)
        --
        ([xshift=-7mm,yshift=-13mm]frame.north east);
    }
]
        \begin{itemize}
    \fi
}
{%
    \end{itemize}
    \ifdefined\HCode
        \HCode{%
</div>
</section>
        }%
    \else
        \end{tcolorbox}
    \fi
}



\addPrintStyle{../../..}

\title{Oefeningen — De regel van l'Hôpital}
\author{Ingmar Herreman}

\begin{document}

% FVD_CHAPTER_DATA_START
% Automatisch gegenereerd uit index.tex.
% Niet handmatig aanpassen.
\fvdchapterdata{3}{Afleiden van samengestelde functie}{3}
% FVD_CHAPTER_DATA_END

\maketitle
\FVDbeginExercises


% ============================================================
% BASIS
% ============================================================



\begin{oefening}[Herken de vorm]{\oefB\ESS}

Bepaal telkens de vorm die ontstaat door substitutie. Beslis of
l'Hôpital rechtstreeks mag worden toegepast.

\begin{question}
\[\lim_{x\to3}\frac{x^2-4x+3}{x-3}.\]

\begin{oplossing}
Substitutie geeft \(0/0\). L'Hôpital mag rechtstreeks worden toegepast.
\end{oplossing}
\end{question}

\begin{question}
\[\lim_{x\to+\infty}\frac{3x-2}{9x+7}.\]

\begin{oplossing}
De teller en noemer worden beide in grootte onbegrensd. De vorm is
\(\infty/\infty\). L'Hôpital mag rechtstreeks worden toegepast.
\end{oplossing}
\end{question}

\begin{question}
\[\lim_{x\to0}\frac{x+1}{x^2+1}.\]

\begin{oplossing}
Substitutie geeft \(1/1\). Er is geen onbepaalde vorm. L'Hôpital mag
niet worden toegepast en de limiet is \(1\).
\end{oplossing}
\end{question}

\begin{question}
\[\lim_{x\to2}\frac1{x-2}.\]

\begin{oplossing}
Dit is geen vorm \(0/0\) of \(\infty/\infty\). L'Hôpital mag niet
rechtstreeks worden toegepast. De eenzijdige limieten moeten worden
onderzocht.
\end{oplossing}
\end{question}

\begin{question}
\[\lim_{x\to+\infty}\left(\sqrt{x^2+1}-x\right).\]

\begin{oplossing}
De vorm is \(\infty-\infty\). L'Hôpital mag niet rechtstreeks worden
toegepast. Rationaliseren is een natuurlijke eerste stap.
\end{oplossing}
\end{question}
\end{oefening}

\begin{oefening}[De vorm \(0/0\)]{\oefB\ESS}

Bereken met de regel van l'Hôpital.

\begin{question}
\[\lim_{x\to3}\frac{x^2-4x+3}{x-3}.\]

\begin{oplossing}
De vorm is \(0/0\). Dus
\[\lim_{x\to3}\frac{x^2-4x+3}{x-3}\overset{\mathrm H}{=}\lim_{x\to3}(2x-4)=2.\]
\end{oplossing}
\end{question}

\begin{question}
\[\lim_{x\to-1}\frac{x+1}{x^2-1}.\]

\begin{oplossing}
De vorm is \(0/0\). Dus
\[\lim_{x\to-1}\frac{x+1}{x^2-1}\overset{\mathrm H}{=}\lim_{x\to-1}\frac1{2x}=-\frac12.\]
\end{oplossing}
\end{question}

\begin{question}
\[\lim_{x\to1}\frac{2x^2+3x-5}{3x^2+4x-7}.\]

\begin{oplossing}
De vorm is \(0/0\). Dus
\[\lim_{x\to1}\frac{2x^2+3x-5}{3x^2+4x-7}\overset{\mathrm H}{=}\lim_{x\to1}\frac{4x+3}{6x+4}=\frac7{10}.\]
\end{oplossing}
\end{question}

\begin{question}
\[\lim_{x\to\sqrt2}\frac{x^4-4}{x^2-2}.\]

\begin{oplossing}
De vorm is \(0/0\). Dus
\[\lim_{x\to\sqrt2}\frac{x^4-4}{x^2-2}\overset{\mathrm H}{=}\lim_{x\to\sqrt2}\frac{4x^3}{2x}=\lim_{x\to\sqrt2}2x^2=4.\]
\end{oplossing}
\end{question}

\begin{question}
\[\lim_{x\to4}\frac{\sqrt{x}-2}{x-4}.\]

\begin{oplossing}
De vorm is \(0/0\). Dus
\[\lim_{x\to4}\frac{\sqrt{x}-2}{x-4}\overset{\mathrm H}{=}\lim_{x\to4}\frac1{2\sqrt{x}}=\frac14.\]
\end{oplossing}
\end{question}

\begin{question}
\[\lim_{x\to9}\frac{\sqrt{x}-3}{x-9}.\]

\begin{oplossing}
De vorm is \(0/0\). Dus
\[\lim_{x\to9}\frac{\sqrt{x}-3}{x-9}\overset{\mathrm H}{=}\lim_{x\to9}\frac1{2\sqrt{x}}=\frac16.\]
\end{oplossing}
\end{question}
\end{oefening}

\begin{oefening}[De vorm \(\infty/\infty\)]{\oefB\ESS}

Bereken met de regel van l'Hôpital. In deze oefening is l'Hôpital
verplicht, ook wanneer je een snellere algebraïsche methode kent.

\begin{question}
\[\lim_{x\to+\infty}\frac{3x-2}{9x+7}.\]

\begin{oplossing}
De vorm is \(\infty/\infty\). Dus
\[\lim_{x\to+\infty}\frac{3x-2}{9x+7}\overset{\mathrm H}{=}\lim_{x\to+\infty}\frac39=\frac13.\]
\end{oplossing}
\end{question}

\begin{question}
\[\lim_{x\to+\infty}\frac{6x^2+2x+2}{5x^2-3x+4}.\]

\begin{oplossing}
Na de eerste toepassing ontstaat opnieuw \(\infty/\infty\):
\[\lim_{x\to+\infty}\frac{6x^2+2x+2}{5x^2-3x+4}\overset{\mathrm H}{=}\lim_{x\to+\infty}\frac{12x+2}{10x-3}\overset{\mathrm H}{=}\frac{12}{10}=\frac65.\]
\end{oplossing}
\end{question}

\begin{question}
\[\lim_{x\to-\infty}\frac{x}{x^2+5}.\]

\begin{oplossing}
De vorm is \((-\infty)/(+\infty)\), dus van het type
\(\infty/\infty\). Daarom
\[\lim_{x\to-\infty}\frac{x}{x^2+5}\overset{\mathrm H}{=}\lim_{x\to-\infty}\frac1{2x}=0.\]
\end{oplossing}
\end{question}

\begin{question}
\[\lim_{x\to+\infty}\frac{2x^3-x+1}{x^2+4}.\]

\begin{oplossing}
De vorm is \(\infty/\infty\). Dus
\[\lim_{x\to+\infty}\frac{2x^3-x+1}{x^2+4}\overset{\mathrm H}{=}\lim_{x\to+\infty}\frac{6x^2-1}{2x}=+\infty.\]
\end{oplossing}
\end{question}
\end{oefening}

\begin{oefening}[Meermaals l'Hôpital]{\oefB\ESS}

Bereken. Controleer na elke toepassing opnieuw de vorm.

\begin{question}
\[\lim_{x\to1}\frac{x^3-3x+2}{(x-1)^2}.\]

\begin{oplossing}
Eerst ontstaat \(0/0\):
\[\lim_{x\to1}\frac{x^3-3x+2}{(x-1)^2}\overset{\mathrm H}{=}\lim_{x\to1}\frac{3x^2-3}{2(x-1)}.\]
Opnieuw ontstaat \(0/0\):
\[\lim_{x\to1}\frac{3x^2-3}{2(x-1)}\overset{\mathrm H}{=}\lim_{x\to1}\frac{6x}{2}=3.\]
\end{oplossing}
\end{question}

\begin{question}
\[\lim_{x\to2}\frac{x^3-6x^2+12x-8}{x^2-4x+4}.\]

\begin{oplossing}
De vorm is \(0/0\). Een eerste toepassing geeft opnieuw \(0/0\):
\[\lim_{x\to2}\frac{x^3-6x^2+12x-8}{x^2-4x+4}\overset{\mathrm H}{=}\lim_{x\to2}\frac{3x^2-12x+12}{2x-4}.\]
Nogmaals l'Hôpital:
\[\lim_{x\to2}\frac{3x^2-12x+12}{2x-4}\overset{\mathrm H}{=}\lim_{x\to2}\frac{6x-12}{2}=0.\]
\end{oplossing}
\end{question}

\begin{question}
\[\lim_{x\to+\infty}\frac{2x^3+x^2+1}{3x^4-2x}.\]

\begin{oplossing}
Na elke toepassing blijft de vorm \(\infty/\infty\), tot de limiet
rechtstreeks kan worden bepaald:
\[
\begin{aligned}
\lim_{x\to+\infty}\frac{2x^3+x^2+1}{3x^4-2x}
&\overset{\mathrm H}{=}\lim_{x\to+\infty}\frac{6x^2+2x}{12x^3-2}\\
&\overset{\mathrm H}{=}\lim_{x\to+\infty}\frac{12x+2}{36x^2}\\
&\overset{\mathrm H}{=}\lim_{x\to+\infty}\frac{12}{72x}=0.
\end{aligned}
\]
\end{oplossing}
\end{question}
\end{oefening}

% ============================================================
% VERDIEPING
% ============================================================


\begin{oefening}[Welke methode kies je?]{\oefU\ESS}

Bereken elke limiet. Bepaal eerst de vorm en kies vervolgens een
efficiënte methode. Motiveer waarom l'Hôpital wel of niet aangewezen is.

\begin{question}
\[\lim_{x\to2}\frac{x^2-4}{x-2}.\]

\begin{oplossing}
De vorm is \(0/0\), dus l'Hôpital mag. Ontbinden is hier eenvoudiger:
\[\frac{x^2-4}{x-2}=x+2\qquad(x\neq2).\]
Daarom is de limiet
\[\boxed{4}.\]
\end{oplossing}
\end{question}

\begin{question}
\[\lim_{x\to+\infty}\frac{4x^2-3x+1}{2x^2+5}.\]

\begin{oplossing}
De vorm is \(\infty/\infty\). L'Hôpital mag, maar delen door \(x^2\)
is efficiënter:
\[\lim_{x\to+\infty}\frac{4-\frac3x+\frac1{x^2}}{2+\frac5{x^2}}=2.\]
\end{oplossing}
\end{question}

\begin{question}
\[\lim_{x\to4}\frac{\sqrt{x}-2}{x-4}.\]

\begin{oplossing}
De vorm is \(0/0\). Zowel l'Hôpital als rationaliseren is geschikt.
Met l'Hôpital:
\[\lim_{x\to4}\frac{\sqrt{x}-2}{x-4}\overset{\mathrm H}{=}\lim_{x\to4}\frac1{2\sqrt{x}}=\frac14.\]
\end{oplossing}
\end{question}

\begin{question}
\[\lim_{x\to+\infty}\left(\sqrt{x^2+3x}-x\right).\]

\begin{oplossing}
De vorm is \(\infty-\infty\). L'Hôpital mag niet rechtstreeks worden
toegepast. Rationaliseren is de natuurlijke eerste stap:
\[
\sqrt{x^2+3x}-x
=\frac{3x}{\sqrt{x^2+3x}+x}
=\frac3{\sqrt{1+\frac3x}+1}.
\]
Daarom
\[\boxed{\frac32}.\]
\end{oplossing}
\end{question}
\end{oefening}

\begin{oefening}[Vind en herstel de fout]{\oefU\ESS}

Analyseer elke redenering. Geef precies aan waar de fout zit en bereken
daarna de limiet correct.

\begin{question}
Een leerling schrijft
\[\lim_{x\to0}\frac{x+1}{x+2}\overset{\mathrm H}{=}\lim_{x\to0}\frac11=1.\]

\begin{oplossing}
De fout gebeurt vóór het afleiden. Substitutie geeft \(1/2\), niet
\(0/0\) of \(\infty/\infty\). L'Hôpital mag dus niet worden toegepast:
\[\lim_{x\to0}\frac{x+1}{x+2}=\frac12.\]
\end{oplossing}
\end{question}

\begin{question}
Een leerling berekent
\[\lim_{x\to1}\frac{x^2-1}{x-1}\]
en schrijft na het herkennen van \(0/0\):
\[\lim_{x\to1}\frac{x^2-1}{x-1}\overset{\mathrm H}{=}\lim_{x\to1}\frac{2x(x-1)-(x^2-1)}{(x-1)^2}.\]

\begin{oplossing}
De leerling gebruikt de quotiëntregel. Dat is niet de regel van
l'Hôpital. Teller en noemer moeten afzonderlijk worden afgeleid:
\[\lim_{x\to1}\frac{x^2-1}{x-1}\overset{\mathrm H}{=}\lim_{x\to1}\frac{2x}{1}=2.\]
\end{oplossing}
\end{question}

\begin{question}
Een leerling schrijft
\[\lim_{x\to+\infty}\frac{3x^2+1}{x^3-2}\overset{\mathrm H}{=}\lim_{x\to+\infty}\frac{6x}{3x^2}\overset{\mathrm H}{=}\lim_{x\to+\infty}\frac6{6x}\overset{\mathrm H}{=}\frac60=0.\]
Analyseer de laatste stap.

\begin{oplossing}
De eerste twee toepassingen zijn geldig, omdat telkens
\(\infty/\infty\) ontstaat. Na de tweede toepassing staat
\[\lim_{x\to+\infty}\frac1x=0.\]
Die limiet is al bepaald. Er mag niet opnieuw l'Hôpital worden toegepast:
de vorm is dan niet meer \(\infty/\infty\). De conclusie \(0\) is wel
juist, maar de laatste toepassing van l'Hôpital is fout.
\end{oplossing}
\end{question}
\end{oefening}

\begin{oefening}[Eenzijdige limieten]{\oefU\ESS}

Bereken en let zorgvuldig op het teken en de richting van de limiet.

\begin{question}
\[\lim_{x\to6^-}\frac{x^2-8x+12}{x^2-12x+36}.\]

\begin{oplossing}
Substitutie geeft \(0/0\). Met l'Hôpital:
\[\lim_{x\to6^-}\frac{x^2-8x+12}{x^2-12x+36}\overset{\mathrm H}{=}\lim_{x\to6^-}\frac{2x-8}{2x-12}.\]
De teller nadert \(4>0\), terwijl de noemer \(0^-\) nadert. Dus
\[\boxed{-\infty}.\]
\end{oplossing}
\end{question}

\begin{question}
\[\lim_{x\to-1^+}\frac{x^2+2x+1}{x^3+3x^2+3x+1}.\]

\begin{oplossing}
De vorm is \(0/0\). Algebraïsch zien we
\[\frac{(x+1)^2}{(x+1)^3}=\frac1{x+1}\qquad(x\neq-1).\]
Voor \(x\to-1^+\) nadert de noemer \(0^+\). Dus
\[\boxed{+\infty}.\]
L'Hôpital kan eveneens worden gebruikt, maar vereenvoudigen is hier
transparanter.
\end{oplossing}
\end{question}

\begin{question}
\[\lim_{x\to\frac12^-}\frac{2x-1}{4x^2-4x+1}.\]

\begin{oplossing}
Omdat \(4x^2-4x+1=(2x-1)^2\),
\[\frac{2x-1}{4x^2-4x+1}=\frac1{2x-1}\qquad\left(x\neq\frac12\right).\]
Voor \(x\to\frac12^-\) nadert de noemer \(0^-\). Dus
\[\boxed{-\infty}.\]
\end{oplossing}
\end{question}

\begin{question}
\[\lim_{x\to0^+}\frac{\sqrt{x}}{x}.\]

\begin{oplossing}
De vorm is \(0/0\). Voor \(x>0\) mag l'Hôpital worden toegepast:
\[\lim_{x\to0^+}\frac{\sqrt{x}}x\overset{\mathrm H}{=}\lim_{x\to0^+}\frac1{2\sqrt{x}}=+\infty.\]
Dus
\[\boxed{+\infty}.\]
\end{oplossing}
\end{question}
\end{oefening}

\begin{oefening}[Eerst omvormen]{\oefU\ESS}

Bereken. L'Hôpital mag niet als eerste stap worden gebruikt. Geef ook
aan of l'Hôpital na de omvorming nog nodig of zinvol is.

\begin{question}
\[\lim_{x\to+\infty}\left(\sqrt{x^2+2}-x\right).\]

\begin{oplossing}
De vorm is \(\infty-\infty\). Rationaliseren geeft
\[\sqrt{x^2+2}-x=\frac2{\sqrt{x^2+2}+x}.\]
De noemer gaat naar \(+\infty\), dus
\[\boxed{0}.\]
Na de omvorming is l'Hôpital niet nodig.
\end{oplossing}
\end{question}

\begin{question}
\[\lim_{x\to+\infty}\left(\sqrt{x^2+x}-x\right).\]

\begin{oplossing}
De vorm is \(\infty-\infty\). Rationaliseren geeft
\[\sqrt{x^2+x}-x=\frac{x}{\sqrt{x^2+x}+x}.\]
Omdat \(x>0\) voor voldoende grote \(x\), delen we teller en noemer door
\(x\):
\[\frac1{\sqrt{1+\frac1x}+1}\longrightarrow\frac12.\]
Dus
\[\boxed{\frac12}.\]
\end{oplossing}
\end{question}

\begin{question}
\[\lim_{x\to+\infty}\left(\sqrt{4x^2+3x}-2x\right).\]

\begin{oplossing}
De vorm is \(\infty-\infty\). Rationaliseren geeft
\[
\sqrt{4x^2+3x}-2x
=\frac{3x}{\sqrt{4x^2+3x}+2x}
=\frac3{\sqrt{4+\frac3x}+2}.
\]
Daarom
\[\boxed{\frac34}.\]
\end{oplossing}
\end{question}
\end{oefening}

\begin{oefening}[Parameterlimiet]{\oefU\ESS}

Neem \(a>0\) en \(m,n\in\mathbb{N}_0\) met \(m\geq1\) en \(n\geq1\).
Bepaal
\[\lim_{x\to a}\frac{x^m-a^m}{x^n-a^n}.\]

\begin{oplossing}
De vorm is \(0/0\). Met l'Hôpital:
\[
\lim_{x\to a}\frac{x^m-a^m}{x^n-a^n}
\overset{\mathrm H}{=}
\lim_{x\to a}\frac{mx^{m-1}}{nx^{n-1}}
=\boxed{\frac mn a^{m-n}}.
\]
\end{oplossing}
\end{oefening}

% ============================================================
% GEVORDERD
% ============================================================



\begin{oefening}[Bepaal de parameter]{\oefG}

Bepaal \(a\in\mathbb{R}\) zodat
\[\lim_{x\to1}\frac{x^2+ax-(a+1)}{x-1}=5.\]

\begin{oplossing}
Voor elke \(a\) is de teller bij \(x=1\) gelijk aan
\(1+a-(a+1)=0\). De vorm is dus \(0/0\). Met l'Hôpital:
\[
\lim_{x\to1}\frac{x^2+ax-(a+1)}{x-1}
\overset{\mathrm H}{=}
\lim_{x\to1}(2x+a)=2+a.
\]
We eisen \(2+a=5\). Dus
\[\boxed{a=3}.\]
\end{oplossing}
\end{oefening}

\begin{oefening}[Ontwerp zelf een limiet]{\oefG}

Zoek twee verschillende veeltermfuncties \(f\) en \(g\) waarvoor
\[
f(2)=g(2)=0
\qquad\text{en}\qquad
\lim_{x\to2}\frac{f(x)}{g(x)}=3.
\]
Je functies mogen niet beide eerstegraads zijn. Toon met l'Hôpital aan
dat jouw voorbeeld aan de voorwaarden voldoet.

\begin{oplossing}
Er zijn veel mogelijkheden. Bijvoorbeeld
\[f(x)=3(x^2-4)\qquad\text{en}\qquad g(x)=2(x^2+x-6).\]
Dan is \(f(2)=g(2)=0\). Met l'Hôpital:
\[
\lim_{x\to2}\frac{3(x^2-4)}{2(x^2+x-6)}
\overset{\mathrm H}{=}
\lim_{x\to2}\frac{6x}{4x+2}
=\frac{12}{10}.
\]
Dit voorbeeld voldoet dus niet: de limiet is \(6/5\), niet \(3\).

We moeten onze keuze aanpassen. Neem bijvoorbeeld
\[f(x)=3(x^2-4)\qquad\text{en}\qquad g(x)=x^2-4.\]
Dan zijn beide functies niet eerstegraads en
\[
\lim_{x\to2}\frac{f(x)}{g(x)}
\overset{\mathrm H}{=}
\lim_{x\to2}\frac{6x}{2x}=3.
\]
Dus dit is een geldig voorbeeld. Andere antwoorden zijn mogelijk.
\end{oplossing}
\end{oefening}

\begin{oefening}[Analyseer een bewering]{\oefG}

Een leerling beweert:

\begin{center}
``Als teller en noemer beide naar oneindig gaan, mag je teller en
noemer blijven afleiden tot er constanten overblijven.''
\end{center}

Beoordeel deze uitspraak. Geef aan wat er correct aan is en wat ontbreekt.

\begin{oplossing}
De uitspraak is te sterk. Bij een vorm \(\infty/\infty\) kan l'Hôpital
worden toegepast wanneer de voorwaarden van de regel vervuld zijn. Na
elke toepassing moet opnieuw worden nagegaan welke vorm ontstaat.

Zodra de nieuwe limiet geen \(0/0\) of \(\infty/\infty\) meer is, mag
l'Hôpital niet automatisch opnieuw worden toegepast. Bovendien hoeft men
niet ``tot constanten'' af te leiden: men stopt zodra de nieuwe limiet
kan worden bepaald.

Ten slotte kan een andere methode, bijvoorbeeld delen door de hoogste
macht, veel efficiënter zijn.
\end{oplossing}
\end{oefening}

\begin{oefening}[Eén probleem, drie methoden]{\oefG}

Beschouw
\[\lim_{x\to2}\frac{x^3-8}{x^2-4}.\]

\begin{question}
Bereken de limiet door te ontbinden.

\begin{oplossing}
\[
\frac{x^3-8}{x^2-4}
=\frac{(x-2)(x^2+2x+4)}{(x-2)(x+2)}
=\frac{x^2+2x+4}{x+2}.
\]
Dus
\[\boxed{3}.\]
\end{oplossing}
\end{question}

\begin{question}
Bereken dezelfde limiet met l'Hôpital.

\begin{oplossing}
De vorm is \(0/0\). Daarom
\[\lim_{x\to2}\frac{x^3-8}{x^2-4}\overset{\mathrm H}{=}\lim_{x\to2}\frac{3x^2}{2x}=3.\]
\end{oplossing}
\end{question}

\begin{question}
Vergelijk beide methoden. Welke informatie moet je controleren vóór je
l'Hôpital gebruikt?

\begin{oplossing}
Ontbinden is hier zeer transparant omdat de gemeenschappelijke factor
\(x-2\) zichtbaar wordt. L'Hôpital is eveneens kort.

Vóór l'Hôpital wordt toegepast, moet eerst worden vastgesteld dat de
limiet de onbepaalde vorm \(0/0\) heeft en dat aan de voorwaarden van de
regel is voldaan.
\end{oplossing}
\end{question}
\end{oefening}



\FVDendExercises
\end{document}


